\documentclass[10pt,a4paper]{amsart}
\usepackage[margin=19mm]{geometry}
\usepackage{amsmath,amssymb,mathtools}
\usepackage[scr=rsfs,cal=euler]{mathalfa}
\usepackage{xfrac}
\usepackage[unicode,hidelinks,bookmarks=false]{hyperref}
\newcommand{\ie}{i.e.\ }
\newcommand{\cf}{cf.\ }
\newcommand{\resp}{respectively\ }
\newcommand{\Subsection}{\subsection}
\providecommand{\nobreakdash}{\nobreak\hbox{-}}
\setlength{\emergencystretch}{2em}
\multlinegap0pt
\usepackage{bm}
\usepackage{bbm}
\usepackage{dsfont}
\usepackage{xspace}
\usepackage{cancel}
\usepackage{mathtools}
\usepackage{amsthm}
\makeatletter
\@ifundefined{thm}{\newtheorem{thm}{Thm}}{}
\@ifundefined{lemma}{\newtheorem{lemma}[thm]{Lemma}}{}
\makeatother
\DeclareMathOperator{\pr}{pr}
\title[Non-squeezing and other global rigidity results in locally c]{Non-squeezing and other global rigidity results in locally conformal symplectic geometry}
\author{M\'elanie Bertelson, Pranav Chakravarthy, Sheila Sandon}
\date{Source: arXiv:2511.16329v2 (CC BY 4.0)}
\begin{document}
\maketitle

\noindent\emph{Source and licence.} Non-squeezing and other global rigidity results in locally conformal symplectic geometry. Version arXiv:2511.16329v2 (CC BY 4.0), distributed under the Creative Commons Attribution 4.0 International licence, as recorded in the arXiv licence metadata for this exact version and shown on the abstract page https://arxiv.org/abs/2511.16329v2. Full rights notice: licenses/arxiv-2511.16329-CC-BY-4.0.txt.\par\smallskip

\noindent\textit{Editorial scope.} Counted unit: the Lemma whose source label is \texttt{lemma: bijection critical points F} in the article (source0.tex lines 4942--4953, source label \texttt{lemma: bijection critical points F}), delivered together with the article's own complete original proof of it (source0.tex lines 4955--5098, one \texttt{proof} environment of 144 lines). No mathematics is written, repaired or re-derived: statement and proof are the article's own text line for line. Required context retained uncounted: equation: pullback of lambda (lines 4924--4927). Every definition, display and label of those lines is retained unchanged. Omitted from this package and not redistributed: 1--4923, 4928--4941, 4954--4954, 5099--8274. Nothing inside a retained excerpt is omitted or altered: every retained body is delivered exactly as the article's lines are written, so no omission receipt of any kind accompanies this package. Citations: the retained excerpts carry no citation command at all, so no bibliography entry of the article is transcribed and no reference is redistributed. Reference adaptation: none was needed and none was made; every reference inside the delivered unit resolves inside it, so no unresolved reference or citation is produced. Printed slips kept literal: no printed word, symbol, hypothesis or number of the counted statement or of its proof was altered. Layout and class: the article's own class is \texttt{amsart} with packages including import, graphicx, amssymb, amsthm, amstext, amsmath, fullpage, tikz-cd, hyperref, tikz, pgfplots; the wrapper is the lane's pinned \texttt{amsart} editorial wrapper at 10pt on A4, which restates the theorem-like environments the retained text needs and adds the article's own macro definitions that the retained text needs (\texttt{\textbackslash pr}), with extra wrapper packages (bm, bbm, dsfont, xspace, cancel, mathtools). The delivered numbering is the unit's own. Assets: no retained span contains a figure environment, a graphics-inclusion command, a PDF-page inclusion, a table or a TikZ picture; no image file is copied into this package, the package declares no asset dependency and no image file is redistributed with it. Licence: version arXiv:2511.16329v2 of this article is distributed under the Creative Commons Attribution 4.0 International licence, as recorded in the arXiv licence metadata for this exact version and shown on the abstract page https://arxiv.org/abs/2511.16329v2. A full rights notice is delivered at licenses/arxiv-2511.16329-CC-BY-4.0.txt. Characters: every retained excerpt is pure ASCII, so the delivered engine, XeLaTeX with its default Latin Modern text font, reports no missing character.
\begin{equation}\label{equation: pullback of lambda}
i_{F, \beta}^{\ast} \, \lambda
= d_{\beta} \big( \left. F \right\lvert_{\Sigma_F} \big) \,.
\end{equation}

\begin{lemma}\label{lemma: bijection critical points F}
Let $F: E \rightarrow \mathbb{R}$ be a generating function
for a Lagrangian submanifold $L$ of $T_{\beta}^{\ast}B$.
Then there is a bijection between the critical points of $F$
and the essential Lee chords from $L$ to the zero section.
The action (hence the time-shift) of any such Lee chord
is equal to the critical value
of the corresponding critical point of $F$.
Moreover,
the bijection sends non-degenerate critical points
to transverse Lee chords.
\end{lemma}

\begin{proof}
By definition,
$L$ is the image of the exact Lagrangian embedding
$i_{F, \beta}: \Sigma_F \rightarrow T^{\ast}_{\beta}B$.
By \eqref{equation: pullback of lambda},
it has action
\[
S: L \rightarrow \mathbb{R} \,,\;
S \big( i_{F, \beta} (e) \big) = F (e) \,.
\]
Suppose that $e = (q, \zeta) \in \Sigma_F$ is a critical point of $F$.
Then
\[
i_{F, \beta} (e) = - F(e) \, \beta (q) \,,
\]
thus $\big( F(e), \gamma_e \big)$ with
\[
\gamma_e: [0, 1] \rightarrow T_{\beta}^{\ast}B \,,\;
\gamma_e(t) = - F(e) \, \beta (q) + t \, F(e) \, \beta (q)
\]
is a Lee chord from the point $- F(e) \, \beta (q)$ of $L$
to the point $0_q$ of the zero section.
The action
\[
S \big(\gamma(0)\big) - S \big(\gamma(1)\big) + \int_{\gamma_e} \lambda
= S \big(\gamma(0)\big) = F(e)
\]
of the Lee chord $\big( F(e), \gamma_e \big)$
is equal to the time-shift,
and to the critical value of the critical point $e$.
In particular,
the Lee chord $\big( F(e), \gamma_e \big)$ is essential.
The map
\begin{equation}\label{equation: bijection critical point F}
e \mapsto \big( F(e), \gamma_e \big)    
\end{equation}
from the set of critical points of $F$
to the set of essential Lee chords
from $L$ to the zero section
is injective,
because $i_{F, \beta}$ is an embedding.
We now show that it is also surjective.
Suppose that $(T, \gamma)$ with
\[
\gamma: [0, 1] \mapsto T_{\beta}^{\ast}B \,,\;
\gamma (t) = \sigma_q + t \, T \, \beta (q)
\]
is an essential Lee chord from $L$ to the zero section.
Since $\sigma_q \in L$,
there is $e = (q, \zeta) \in \Sigma_F$
with $i_{F, \beta} (e) = \sigma_q$.
Thus
\[
dF (e) (\hat{X})
= \sigma_q (X) + F(e) \beta(q) (X)
\]
for all $X \in T_qB$ and $\hat{X} \in T_eE$ projecting to $X$.
Since $(T, \gamma)$ is essential,
the time-shift $T$
is equal to the action $F(e)$.
Since moreover $\gamma(1) = \sigma_q + T \beta(q)$ belongs to the zero section,
we thus have $dF(e) (\hat{X}) = 0$,
and so $e$ is a critical point of $F$.
Moreover, $\sigma_q = - T \beta(q) = - F(e) \beta(q)$
and so $(T, \gamma) = \big( F(e), \gamma_e \big)$.
This shows that the map \eqref{equation: bijection critical point F}
is surjective,
hence a bijection.

Finally,
we show that a critical point $e$ is non-degenerate
if and only if the Lee chord $\big( F(e), \gamma_e \big)$ is transverse.
By definition,
the Lee chord $\big( F(e), \gamma_e \big)$ is not transverse
if and only if
\[
(\varphi_{F(e)}^{\omega})_{\ast} \big(T_{i_{F, \beta}(e)} L \big)
+ T_{0_q} 0_B \neq T_{0_q} \, T^{\ast}B \,,
\]
thus if and only if
$(\varphi_{F(e)}^{\omega})_{\ast} \big(T_{i_{F, \beta}(e)} L \big)
\cap T_{0_q} 0_B \neq \{0\}$,
and so
if and only if there is $Y \in T_e \Sigma_F$
such that
\begin{equation}\label{equation: transversality}
(\varphi_{F(e)}^{\omega} \circ i_{F, \beta})_{\ast} \, (Y) \in T_{0_q} 0_B \,.
\end{equation}
Write $Y = \dot{\delta}(t)$
for a path $\delta: (-\varepsilon, \varepsilon) \to \Sigma_F$
with $\delta (0) = e$.
Set $\delta_o = \pr_1 \circ \,\delta$,
where $\pr_1 : \Sigma_F \subset E = B \times \mathbb{R}^N \to B$
is the projection to the first factor.
Observe that $N_E^*$ is canonically isomorphic
to the pullback bundle $\pr_1^{\,\ast}\, T^*B$,
and denote by $p: N_E^* \to T^*B$ the natural projection.
We can then write $i_{F,\beta} (e) = p \, (d_\beta F(e))$ and
\begin{align*}
(\varphi_{F(e)}^{\omega} \circ i_{F, \beta})_{\ast} \, (Y)
&=  \left. \frac{d}{dt} \right \lvert_{t=0}
p \, \Big( d_\beta F \big(\delta(t)\big) \Big) + F(e) \, \beta \big(\delta_o(t)\big) \\
&= \left.\frac{d}{dt} \right \lvert_{t=0}
p \, \Big( dF \big(\delta(t)\big) \Big)
+ \Big(F(e) - F \big(\delta(t)\big)\Big) \, \beta \big(\delta_o(t)\big) \,.
\end{align*}
We observe that $\delta_1 (t) := p \, \left(dF(\delta(t))\right)$
and $\delta_2 (t) := \big( F(e) - F(\delta(t)) \big) \, \beta(\delta_o(t))$
are two paths in $T^*B$ with $\delta_1 (0) = \delta_2(0) = 0_q$.
Therefore,
their tangent vectors at $t = 0$
both belong to $T_{0_q} T^*B$.
The latter vector space admits a canonical decomposition
as the direct sum of $T_q B$ -- the ``horizontal" part -- and $T^*_qB$ -- the ``vertical" part.
The decompositions of $\left. \frac{d}{dt} \right\lvert_{t=0} \delta_1(t)$
and $\left. \frac{d}{dt} \right\lvert_{t=0} \delta_2(t)$
into horizontal and vertical parts are
$$
\left.\frac{d}{dt}\right\lvert_{t=0} \delta_1 (t)
= p_{*} \big( Y + d^2F (Y, \,\cdot\,) \big)
= {\pr_1}_{*} (Y) + p_{\ast} \big( d^2F(Y, \,\cdot\,) \big)
$$
and
$$
\left.\frac{d}{dt}\right\lvert_{t=0} \delta_2 (t)
= {\pr_1}_{*} (Y) + \left.\frac{d}{dt}\right\lvert_{t=0}
\Big( F(e) - F \big(\delta(q)\big) \Big) \, \beta(q)
= {\pr_1}_{*} (Y) + 0_q \,,
$$
where we have used that $e$ is a critical point of $F$.
The condition \eqref{equation: transversality}
is equivalent to asking that the sum of the vertical parts
of $\left. \frac{d}{dt} \right\lvert_{t=0} \delta_1(t)$
and $\left. \frac{d}{dt} \right\lvert_{t=0} \delta_2(t)$
vanishes.
Since this sum
is $p_{\ast} \big(d^2F (Y, \,\cdot\,)\big)$,
we conclude that 
the Lee chord $\big( F(e), \gamma_e \big)$ is not transverse
if and only if there is $Y \in T_e\Sigma_F$
such that $p_{\ast} \big(d^2F (Y, \,\cdot\,)\big)$ vanishes,
hence $d^2F (Y, \,\cdot\,)$ vanishes,
and so if and only if $e$ is a degenerate critical point of $F$.
\end{proof}

\end{document}
